% Beamer presentation — Adaptive difficulty: Lifetime & Workspace modes
% Paste into Overleaf. Compile with pdfLaTeX.

\documentclass[10pt,aspectratio=169]{beamer}

\usetheme{Madrid}
\usecolortheme{seahorse}
\setbeamertemplate{navigation symbols}{}
\setbeamertemplate{footline}[frame number]

\usepackage{amsmath, amssymb}
\usepackage{booktabs}
\usepackage{xcolor}

\newcommand{\target}{r}
\newcommand{\obs}{s_n}
\newcommand{\err}{e_n}
\newcommand{\Kp}{K_p}
\newcommand{\Ki}{K_i}
\newcommand{\Iaccum}{I_n}
\newcommand{\dband}{\delta}

\title[Adaptive Difficulty]{Adaptive Difficulty }
\subtitle{Lifetime and Workspace Modes}
\author{}
\institute{}
\date{}

\begin{document}

\frame{\titlepage}

% ============================================================
\begin{frame}{System overview}
\begin{itemize}
    \item Goal: keep each patient's success rate at a daily target $\target_d$, drawn each day from a Gaussian centred on the patient's assigned group mean.
    \item Three patient groups in the study: group means $\mu \in \{0.7,\ 0.8,\ 0.9\}$, with standard deviation $\sigma = 0.05$ across days.
    \item 14 daily targets per patient are sampled in advance: $\target_d \sim \mathcal{N}(\mu, \sigma^2)$ for $d = 1, \ldots, 14$.
    \item The controller treats each day independently: receives the day's $\target_d$, drives the observed success rate toward it.
    \item Two interchangeable difficulty modes:
    \begin{itemize}
        \item \textbf{Lifetime mode} --- difficulty $=$ apple lifetime (seconds).
        \item \textbf{Workspace mode} --- difficulty $=$ spawn distance from workspace centre.
    \end{itemize}
    \item Both use the same calibration and controller method; only the units differ.
\end{itemize}
\end{frame}

% ============================================================
\begin{frame}{Assigned target rate --- per-day Gaussian sampling}
\textbf{Per-patient setup (Day 0):}
\begin{enumerate}
    \item Patient is randomised to one of three groups: $\mu \in \{0.70,\ 0.80,\ 0.90\}$.
    \item Standard deviation fixed at $\sigma = 0.05$ for all groups (matches the typical SD of session-to-session variation in pilot data).
    \item Pre-sample $14$ daily targets:
    \[
    \target_d \sim \mathcal{N}(\mu, \sigma^2), \qquad d = 1, \ldots, 14
    \]

\end{enumerate}

\end{frame}

% ============================================================
\begin{frame}{Trial structure (both modes)}
\begin{itemize}
    \item \textbf{Trial 1 --- calibration.} Apple difficulty is randomised across the full range. Catch/miss outcomes are logged.
    \item At end of Trial 1, estimate the patient's $50\%$ catch \textbf{threshold} $t$.
    \item \textbf{Trial 2 onward --- control.} Each apple's difficulty is sampled from a window of width $W$ centred at $t + \Delta_n$.
    \item After each trial, the PI controller updates $\Delta$ based on observed success rate.
\end{itemize}
\end{frame}

% ============================================================
\section{Lifetime mode}

\begin{frame}{Lifetime mode --- variables}
\begin{center}
\begin{tabular}{lll}
\toprule
Symbol & Meaning & Units \\
\midrule
$\ell$ & apple lifetime & s \\
$t$ & threshold ($50\%$ catch lifetime) & s \\
$\Delta_n$ & controller offset after trial $n$ & s \\
$W$ & window width (fixed) & s \\
$\target$ & assigned target success rate (this day) & --- \\
$\obs$ & observed catch rate in trial $n$ & --- \\
$\err$ & error signal $= \obs - \target$ & --- \\
$\Iaccum$ & integral accumulator (sum of errors) & --- \\
$u_n$ & correction applied to offset after trial $n$ & s \\
$\Kp$ & proportional gain & --- \\
$\Ki$ & integral gain & --- \\
$\dband$ & dead-band half-width & --- \\
\bottomrule
\end{tabular}
\end{center}
\vspace{0.5em}
Fixed constants: $W = 2.4$ s, $\Kp = 0.35$, $\Ki = 0.05$, $\dband = 0.05$, $\ell_{\min} = 0.1$ s, $\ell_{\max} = 8$ s.
\end{frame}

% ============================================================
\begin{frame}{Lifetime mode --- Trial 1 calibration}
\textbf{Per-apple lifetime (random across full range):}
\[
\ell \sim \mathcal{U}(\ell_{\min},\, \ell_{\max})
\]

\textbf{At end of Trial 1, from the catch/miss log:}
\begin{align*}
\ell_{\min}^{\text{caught}} &= \min \{\ell : \text{apple was caught}\} \\
\ell_{\max}^{\text{missed}} &= \max \{\ell : \text{apple was missed}\}
\end{align*}

\textbf{Threshold (midpoint between extremes --- estimator of the $50\%$ catch point):}
\[
\boxed{t = \tfrac{1}{2}\bigl(\ell_{\min}^{\text{caught}} + \ell_{\max}^{\text{missed}}\bigr)}
\]

\textbf{Initial controller offset (set so the sampling window is biased toward the target rate):}
\[
\Delta_0 = (\target - 0.5)\, W
\]
\end{frame}

% ============================================================
\begin{frame}{Lifetime mode --- sampling in Trial 2 onward}
\textbf{Window centre:}
\[
c_n = t + \Delta_n
\]

\textbf{Per-apple lifetime (uniform within the window):}
\[
\ell \sim \mathcal{U}\!\left( c_n - \tfrac{W}{2},\ c_n + \tfrac{W}{2} \right)
\]
clamped to $[\ell_{\min}, \ell_{\max}]$.

\vspace{1em}
\textbf{Geometric interpretation of the initial offset:}
\begin{itemize}
    \item Window spans from $t - (1-\target)W$ to $t + \target W$.
\end{itemize}
\end{frame}

% ============================================================
\begin{frame}{Lifetime mode --- PI controller update}
\textbf{After trial $n$:}
\begin{align*}
\obs &= \frac{\text{caught}_n}{\text{spawned}_n}, \qquad \err = \obs - \target \\[3pt]
\Iaccum &= I_{n-1} + \err \qquad \text{(integral accumulator updated every trial)}
\end{align*}

\textbf{Correction (integral runs always, proportional only outside dead band):}
\[
u_n =
\begin{cases}
\Kp \cdot \err + \Ki \cdot \Iaccum & \text{if } |\err| > \dband \\[4pt]
\Ki \cdot \Iaccum & \text{otherwise}
\end{cases}
\]

\textbf{Offset update:}
\[
\Delta_{n+1} = \mathrm{clip}\bigl(\Delta_n - u_n,\ \Delta_{\min},\ \Delta_{\max}\bigr)
\]

\textbf{Sign convention:}
\begin{itemize}
    \item $\err > 0$ $\Rightarrow$ observed rate exceeds target $\Rightarrow$ $u_n > 0$ $\Rightarrow$ $\Delta$ decreases $\Rightarrow$ window shifts toward harder difficulties.
    \item $\err < 0$ $\Rightarrow$ observed rate below target $\Rightarrow$ $u_n < 0$ $\Rightarrow$ $\Delta$ increases $\Rightarrow$ window shifts toward easier difficulties.
\end{itemize}

\textbf{Bounds:} $\Delta_{\min} = \ell_{\min} - t + \tfrac{W}{2}, \quad \Delta_{\max} = \ell_{\max} - t - \tfrac{W}{2}$
\end{frame}

% ============================================================
\begin{frame}{Why PI and not just P?}
\textbf{Limitation of pure P:}
\begin{itemize}
    \item Inside the dead band ($|\err| \le \dband$), the proportional term contributes zero.
    \item If the observed rate settles at a small persistent offset from target (e.g.\ $77\%$ vs target $80\%$), the error lies inside the dead band, no correction is applied, and the offset is never reduced.
\end{itemize}

\textbf{Role of the integral term:}
\begin{itemize}
    \item $\Iaccum$ accumulates all past errors.
    \item $\Ki \cdot \Iaccum$ is applied at every trial, including those inside the dead band.
    \item Provides slow continuous correction when the error is small, and an additive correction when the error is large.
    \item Eliminates the steady-state offset that pure proportional control leaves behind.
\end{itemize}

\textbf{Limitations:}
\begin{itemize}
    \item Introduces slower-timescale dynamics in the closed-loop response.
    \item Susceptible to integral windup if $\Delta$ saturates at its bounds; no anti-windup mechanism is implemented in this version.
\end{itemize}
\end{frame}

% ============================================================
\section{Workspace mode}

\begin{frame}{Workspace mode --- the rectangle}
During Trial 1, every cursor position $(x, y)$ is observed.

\vspace{0.5em}
\textbf{Bounding rectangle:}
\begin{align*}
\mathbf{ws}_{\min} &= \bigl(\min_i x_i,\ \min_i y_i\bigr) \\
\mathbf{ws}_{\max} &= \bigl(\max_i x_i,\ \max_i y_i\bigr)
\end{align*}

\textbf{Derived geometry:}
\begin{align*}
\mathbf{ws}_c &= \tfrac{1}{2}(\mathbf{ws}_{\min} + \mathbf{ws}_{\max}) \quad \text{(centre)} \\
d_x &= \tfrac{1}{2}(\mathbf{ws}_{\max,x} - \mathbf{ws}_{\min,x}) \quad \text{(half-width)} \\
d_y &= \tfrac{1}{2}(\mathbf{ws}_{\max,y} - \mathbf{ws}_{\min,y}) \quad \text{(half-height)}
\end{align*}

The rectangle defines the patient's reachable region for this session.
\end{frame}

% ============================================================
\begin{frame}{Workspace mode --- rect-scale}
For an apple at position $\mathbf{p}$, define the displacement vector and angle:
\[
\mathbf{v} = \mathbf{p} - \mathbf{ws}_c, \quad \theta = \mathrm{atan2}(v_y,\, v_x)
\]

\textbf{Distance from centre to rectangle edge in direction $\theta$:}
\[
r_{\text{edge}}(\theta) = \min\!\left( \frac{d_x}{|\cos\theta|},\ \frac{d_y}{|\sin\theta|} \right)
\]

\textbf{Rect-scale (dimensionless difficulty, $\rho \in [0, 1]$):}
\[
\boxed{\rho(\mathbf{p}) = \frac{\|\mathbf{v}\|}{r_{\text{edge}}(\theta)}}
\]

\begin{itemize}
    \item $\rho = 0$ $\Rightarrow$ apple at centre (easy).
    \item $\rho = 1$ $\Rightarrow$ apple at edge (hard).
    \item $\rho > 1$ $\Rightarrow$ apple outside reachable rectangle.
\end{itemize}
\end{frame}

% ============================================================
\begin{frame}{Workspace mode --- variables}
\begin{center}
\begin{tabular}{lll}
\toprule
Symbol & Meaning & Units \\
\midrule
$\rho$ & rect-scale of an apple & dimensionless, $[0,1]$ \\
$t_w$ & workspace threshold & dimensionless \\
$\Delta_n^w$ & controller offset after trial $n$ & dimensionless \\
$W_w$ & window width fraction & dimensionless \\
$\target,\obs,\err,\Iaccum,\Kp,\Ki,\dband$ & (same as lifetime mode) & --- \\
\bottomrule
\end{tabular}
\end{center}

\vspace{0.5em}
Fixed constants: $W_w = 0.30$, $\Kp = 0.35$, $\Ki = 0.05$, $\dband = 0.05$.

In workspace mode the apple lifetime is held constant at $\ell_{\max} = 8$ s so that reach time does not bias the difficulty estimate.
\end{frame}

% ============================================================
\begin{frame}{Workspace mode --- Trial 1 calibration}
\textbf{Per-apple position (uniform across viewport):}
\[
(x, y) \sim \mathcal{U}\bigl([0.10\,w,\ 0.85\,w]\bigr) \times \mathcal{U}\bigl([0.10\,h,\ 0.85\,h]\bigr)
\]

\textbf{At end of Trial 1, from the catch/miss log:}
\begin{align*}
\rho_{\max}^{\text{caught}} &= \max \{\rho : \text{apple was caught}\} \\
\rho_{\min}^{\text{missed}} &= \min \{\rho : \text{apple was missed}\}
\end{align*}

\textbf{Threshold (midpoint between extremes):}
\[
\boxed{t_w = \tfrac{1}{2}\bigl(\rho_{\max}^{\text{caught}} + \rho_{\min}^{\text{missed}}\bigr)}
\]

\textbf{Initial controller offset:}
\[
\Delta_0^w = (\target - 0.5)\, W_w
\]
\end{frame}

% ============================================================
\begin{frame}{Workspace mode --- sampling in Trial 2 onward}
\textbf{Window centre (note: minus sign --- positive offset means easier $=$ inward):}
\[
c_n^w = t_w - \Delta_n^w
\]

\textbf{Per-apple rect-scale (uniform within window):}
\[
\rho \sim \mathcal{U}\!\left( c_n^w - \tfrac{W_w}{2},\ c_n^w + \tfrac{W_w}{2} \right)
\]
clamped to $\rho \ge 0$.

\textbf{Per-apple angle (uniform):}
\[
\theta \sim \mathcal{U}(0,\ 2\pi)
\]

\textbf{Spawn position:}
\[
\mathbf{p} = \mathbf{ws}_c + (\cos\theta,\ \sin\theta) \cdot \rho \cdot r_{\text{edge}}(\theta)
\]
\end{frame}

% ============================================================
\begin{frame}{Workspace mode --- PI controller update}
\textbf{After trial $n$:}
\[
\obs = \frac{\text{caught}_n}{\text{spawned}_n}, \qquad \err = \obs - \target, \qquad \Iaccum = I_{n-1} + \err
\]

\textbf{Correction (identical template to lifetime mode):}
\[
u_n =
\begin{cases}
\Kp \cdot \err + \Ki \cdot \Iaccum & \text{if } |\err| > \dband \\[4pt]
\Ki \cdot \Iaccum & \text{otherwise}
\end{cases}
\]

\textbf{Offset update:}
\[
\Delta_{n+1}^w = \mathrm{clip}\bigl(\Delta_n^w - u_n,\ \Delta_{\min}^w,\ \Delta_{\max}^w\bigr)
\]

\textbf{Bounds:}
\[
\Delta_{\min}^w = -t_w + \tfrac{W_w}{2}, \quad \Delta_{\max}^w = (1 - t_w) - \tfrac{W_w}{2}
\]
ensuring the window stays inside $[0, 1]$.
\end{frame}

% ============================================================
\section{Refinements}

\begin{frame}{Two refinements added to the PI system}
After the initial PI controller, two improvements were made:
\begin{enumerate}
    \item \textbf{Staircase calibration}  --- a more accurate replacement for the midpoint-of-extremes threshold estimator used in Trial 1.
    \item \textbf{Reachability constraint}  --- a post-spawn safety clamp preventing apples from appearing further than the patient can physically reach in the chosen lifetime.
\end{enumerate}
The PI controller itself was unchanged in both cases.
\end{frame}

% ============================================================
\begin{frame}{Refinement 1 --- staircase calibration}
\textbf{Motivation:} the midpoint-of-extremes estimator was crude --- one lucky catch or unlucky miss could bias the threshold significantly. A standard 1-up 1-down adaptive staircase concentrates measurements near the $50\%$ point.

\textbf{Procedure:}
\begin{enumerate}
    \item Start at the hardest end of the range: $\ell_0 = \ell_{\min} = 0.1$ s.
    \item For each apple:
    \begin{itemize}
        \item If \textbf{caught} $\Rightarrow$ make harder: $\ell_{n+1} = \ell_n - s$.
        \item If \textbf{missed} $\Rightarrow$ make easier: $\ell_{n+1} = \ell_n + s$.
    \end{itemize}
    \item Track \emph{reversals} (points where the direction changes).
    \item Use a coarse step $s_c = 1.5$ s for the first 2 reversals, then a fine step $s_f = 0.5$ s.
    \item Stop after $N_r = 6$ reversals.
\end{enumerate}

\textbf{Threshold estimate (mean of reversal values):}
\[
t = \frac{1}{N_r} \sum_{i=1}^{N_r} \ell^{(\text{rev}_i)}
\]
\end{frame}

% ============================================================
% \begin{frame}{Refinement 1 --- why staircase is better}
% \begin{itemize}
%     \item \textbf{Concentrates samples near the threshold.} After a few reversals, the staircase bounces around the $50\%$ point. Most of the $\sim$15 trial samples land near the threshold, not spread uniformly over the entire $[\ell_{\min}, \ell_{\max}]$ range.
%     \item \textbf{Self-correcting.} An unlucky miss makes the next apple easier $\Rightarrow$ likely catch $\Rightarrow$ direction reverses $\Rightarrow$ reversal recorded. Random noise is absorbed by the procedure.
%     \item \textbf{Coarse-then-fine schedule.} Large initial step quickly brackets the threshold; smaller step refines the estimate without wasting trials.
%     \item \textbf{Standard psychophysical method.} The 1-up 1-down staircase is known to converge on the $50\%$ point of the psychometric function (Levitt, 1971), giving the procedure a published statistical foundation.
% \end{itemize}

% \textbf{Calibration trial timing change:} Trial 1 timer no longer enforces a strict 60-second cutoff --- calibration ends only when the staircase completes its 6 reversals. Fallback to the old midpoint estimator if the timer fires before staircase completion.
% \end{frame}

% ============================================================
\begin{frame}{Refinement 2 --- reachability constraint}
\textbf{Variables:} $\ell$ apple lifetime (s), $\hat{v}$ estimated speed (px/s), $\mathbf{p}_s$ spawn, $\mathbf{p}_p$ player.

\vspace{0.4em}
\textbf{Speed estimation (running max over catches $i$):}
\[
\hat{v} \leftarrow \max\!\left( \hat{v},\ d_i / \ell_i \right), \qquad d_i = \|\mathbf{p}_s^{(i)} - \mathbf{p}_p^{(i)}\|
\]

\textbf{Maximum reachable distance:}
\[
d_{\max} = \ell \cdot \hat{v}
\]

\textbf{Constraint applied after random spawn:}
Let $\mathbf{u} = \mathbf{p}_s - \mathbf{p}_p$ and $d = \|\mathbf{u}\|$.
\[
\mathbf{p}_s \leftarrow
\begin{cases}
\mathbf{p}_p + (\mathbf{u} / d) \cdot d_{\max} & \text{if } d > d_{\max} \\[4pt]
\mathbf{p}_s & \text{otherwise}
\end{cases}
\]

If the spawn is further than the patient can physically reach in time $\ell$, it is pulled back along the same direction so that distance equals $d_{\max}$.
\end{frame}

% ============================================================
% \begin{frame}{Refinement 2 --- reachability constraint (notes)}
% \textbf{Why ``maximum'' speed estimate (not mean):}
% \begin{itemize}
%     \item Any successful catch demonstrates that the patient \emph{can} cover that distance in that time.
%     \item Maximum across catches is therefore a conservative lower bound on their true peak capability.
%     \item Initialised at $\hat{v} = 0$ (no catches yet); grows monotonically over the session.
% \end{itemize}

% \textbf{Effect on the PI controller:}
% \begin{itemize}
%     \item The controller itself was unchanged --- still PI on offset.
%     \item The reachability constraint was added \emph{after} the spawn position was chosen, as a post-hoc safety clamp.
%     \item Removed the failure mode where the controller's window-sampled lifetime was too short to physically reach the chosen position.
% \end{itemize}
% \end{frame}

% ============================================================
\section{Comparison}

\begin{frame}{Lifetime vs Workspace --- side by side}
\begin{center}
\begin{tabular}{lll}
\toprule
& \textbf{Lifetime} & \textbf{Workspace} \\
\midrule
Difficulty parameter & lifetime $\ell$ (s) & rect-scale $\rho$ \\
Easier direction & longer $\ell$ & smaller $\rho$ (inward) \\
Harder direction & shorter $\ell$ & larger $\rho$ (outward) \\
Threshold & $t$ (s) & $t_w$ \\
Window width & $W = 2.4$ s & $W_w = 0.30$ \\
Window centre & $t + \Delta$ & $t_w - \Delta$ \\
Update rule & $\Delta \leftarrow \Delta - (K_p e + K_i I)$ & $\Delta^w \leftarrow \Delta^w - (K_p e + K_i I)$ \\
\bottomrule
\end{tabular}
\end{center}

\vspace{0.8em}
Identical controller math, identical calibration logic, identical update rule.\\
Only the \textbf{units} (seconds vs dimensionless) and the \textbf{quantity being calibrated} (lifetime vs rect-scale) differ.
\end{frame}

% ============================================================
\section{Fitts' Law system}

\begin{frame}{Motivation --- moving beyond PI on lifetime}
\textbf{Limitations of the PI + staircase system:}
\begin{itemize}
    \item Lifetime alone is not a good difficulty axis. The same lifetime is easy for a close apple and impossible for a far one. The controller did not see this.
    \item One difficulty threshold per patient (from staircase) collapsed all task geometries into a single number, losing information.
    \item No principled model relating task geometry to reach time.
\end{itemize}

\textbf{Fitts' Law provides a principled model:}
\begin{itemize}
    \item Movement time is a function of target distance $A$ and target width $W$.
    \item Two patient-specific parameters $(a, b)$ capture the patient's motor capability.
    \item Difficulty can be set per apple by choosing $A$ and $W$, with lifetime derived from the predicted movement time.
\end{itemize}
\end{frame}

% ============================================================
\begin{frame}{Fitts' Law --- equation and Index of Difficulty}
\textbf{The law (Shannon formulation, MacKenzie 1992):}
\[
\boxed{\text{MT} = a + b \cdot \text{ID}}
\]

\textbf{Index of Difficulty:}
\[
\text{ID} = \log_2\!\left( \frac{A}{W} + 1 \right)
\]
\begin{itemize}
    \item $A$ --- amplitude (distance from current position to target centre), in pixels.
    \item $W$ --- target width (diameter of the catch zone), in pixels.
    \item ID is dimensionless --- larger ID means harder reach (further and/or narrower target).
    \item Typical range in this game: $\text{ID} \in [1, 4]$ bits.
\end{itemize}

\textbf{Parameters:}
\begin{itemize}
    \item $a$ (s) --- intercept; reaction/initiation time independent of difficulty.
    \item $b$ (s/bit) --- slope; additional time per bit of difficulty. Captures motor capability.
\end{itemize}
\end{frame}

% ============================================================
\begin{frame}{Calibration --- Phase 0a (workspace scan)}
\textbf{Goal:} discover which regions of the screen the patient can actually reach with the device, before the Fitts' calibration begins. Replaces the implicit bounding-box workspace of the PI era with an explicit set of reachable cells.

\textbf{Procedure (grid-spiral):}
\begin{enumerate}
    \item Tile the screen with a grid of cells, each $120 \times 120$ px (cell diameter matches the largest Fitts width $W$).
    \item Group cells by ring index $k = \mathrm{round}(\|\mathbf{p}_{\text{cell}} - \mathbf{p}_{\text{centre}}\| / 120)$ (distance from screen centre, quantised by cell size).
    \item Present cells one at a time, inner rings first; within each ring the order is shuffled.
    \item Each cell is presented as a target with fixed lifetime $\ell = 6$ s.
    \item Record hit / miss outcome for each cell.
    \item \textbf{Early termination:} if no cell in a ring is hit, the scan stops --- assume outer rings are unreachable.
\end{enumerate}

\textbf{Outputs (used by Phase 0c):}
\begin{itemize}
    \item \texttt{reachable\_cells} --- the subset of cells the patient hit.
    \item Workspace centre $=$ centroid of reachable cells.
    \item $a_{\text{comfortable}}$ $=$ median distance from workspace centre to reachable cells.
    \item Maximum reach radius $R_{ws}$ $=$ furthest hit cell from screen centre.
\end{itemize}
\end{frame}

% ============================================================
\begin{frame}{Calibration --- Phase 0c}
\textbf{Goal:} measure $(a, b)$ for this patient at the start of every session.

\textbf{Design:}
\begin{itemize}
    \item $5$ distances $A_1, \ldots, A_5$ spanning $0.20 R_{ws}$ to $0.85 R_{ws}$ (where $R_{ws}$ is the patient's workspace radius from Phase 0a).
    \item $3$ widths $W \in \{120, 60, 25\}$ px (large, medium, small).
    \item All $5 \times 3 = 15$ pairs presented, each repeated $5$ times.
    \item Total: $75$ calibration apples.
\end{itemize}

\textbf{Per apple:}
\begin{enumerate}
    \item Pick $(A, W)$ pair.
    \item Spawn target at distance $A$ from current player position, random angle.
    \item Apple radius $= W / 2$.
    \item Record movement time:
    \[
    \text{MT} = t_{\text{caught}} - t_{\text{spawn}} - t_{\text{hold}}
    \]
    where $t_{\text{hold}} = 1$ s is the catch-confirmation hold.
\end{enumerate}
\end{frame}

% ============================================================
\begin{frame}{Calibration --- batch fit of $(a, b)$}
At the end of Phase 0c, fit the model $\text{MT}_k = a + b \cdot \text{ID}_k$ across the 15 pairs by ordinary least squares.

\textbf{Per pair $k$:}
\begin{itemize}
    \item $\text{ID}_k = \log_2(A_k / W_k + 1)$
    \item $\overline{\text{MT}}_k = $ mean movement time across the $5$ valid catches for that pair.
\end{itemize}

\textbf{OLS normal equations:}
\[
b = \frac{N \sum_k \text{ID}_k \, \overline{\text{MT}}_k - \sum_k \text{ID}_k \sum_k \overline{\text{MT}}_k}{N \sum_k \text{ID}_k^2 - (\sum_k \text{ID}_k)^2}
\]
\[
a = \frac{1}{N}\!\left( \sum_k \overline{\text{MT}}_k - b \sum_k \text{ID}_k \right)
\]
with $N = 15$ pairs.

\textbf{Per-pair residual variance (initialised from calibration data):}
\[
\sigma_k^2 = \frac{1}{n_k - 1}\!\left( \sum_i \text{MT}_i^2 - \frac{(\sum_i \text{MT}_i)^2}{n_k} \right)
\]
\end{frame}

% ============================================================
\begin{frame}{Setting apple lifetime in the session}
For each apple, sample an $(A, W)$ pair uniformly at random from the $15$ calibration pairs, then choose lifetime so the patient catches it with probability $r$ (the assigned target rate).

\textbf{Step 1 --- predicted movement time:}
\[
\widehat{\text{MT}} = a + b \cdot \text{ID}, \qquad \text{ID} = \log_2(A/W + 1)
\]

\textbf{Step 2 --- residual standard deviation for this pair:}
\[
\sigma = \sqrt{\sigma_k^2}
\]

\textbf{Step 3 --- lifetime via inverse normal CDF:}
\[
\boxed{\ell = \mathrm{clip}\bigl(\widehat{\text{MT}} + z(r) \cdot \sigma,\ \ell_{\min},\ \ell_{\max}\bigr)}
\]

where $z(r) = \Phi^{-1}(r)$ is the standard-normal quantile for catch rate $r$:
\begin{itemize}
    \item $r = 0.5 \Rightarrow z = 0.00$ (lifetime $= \widehat{\text{MT}}$, $\sim$50\% catch).
    \item $r = 0.7 \Rightarrow z = 0.52$.
    \item $r = 0.8 \Rightarrow z = 0.84$.
    \item $r = 0.9 \Rightarrow z = 1.28$.
\end{itemize}

This assumes movement times for a given $(A, W)$ are approximately normally distributed.
\end{frame}

% ============================================================
\begin{frame}{Online adaptation --- RLS for $(a, b)$}
The patient's $(a, b)$ may drift during the session (warm-up, fatigue). Update them online after every successful catch using \textbf{recursive least squares}.

\textbf{State:}
\[
\boldsymbol{\theta} = \begin{bmatrix} a \\ b \end{bmatrix}, \qquad
\mathbf{P} \in \mathbb{R}^{2 \times 2} \text{ (covariance, initialised as } 1000 \cdot \mathbf{I})
\]

\textbf{Per observation $(\text{MT}, \text{ID})$:}
\[
\mathbf{x} = \begin{bmatrix} 1 \\ \text{ID} \end{bmatrix}, \qquad
e = \text{MT} - \mathbf{x}^\top \boldsymbol{\theta}
\]
\[
\mathbf{g} = \frac{\mathbf{P}\mathbf{x}}{1 + \mathbf{x}^\top \mathbf{P} \mathbf{x}}, \qquad
\boldsymbol{\theta} \leftarrow \boldsymbol{\theta} + \mathbf{g} \cdot e, \qquad
\mathbf{P} \leftarrow \mathbf{P} - \mathbf{g} \mathbf{x}^\top \mathbf{P}
\]

\begin{itemize}
    \item Computationally cheap (no matrix inversion past initialisation; only a $2 \times 2$ update).
    \item Equivalent to OLS on all data seen so far, but updated incrementally.
\end{itemize}
\end{frame}

% ============================================================
\begin{frame}{Online adaptation --- Welford for per-pair $\sigma$}
The per-pair residual SD $\sigma_k$ also updates online, using \textbf{Welford's algorithm} for numerically stable streaming variance.

\textbf{State per pair $k$:} $n_k$ (count), $\mu_k$ (mean), $M_k$ (sum of squared deviations from mean).

\textbf{Per residual $\varepsilon = \text{MT} - (a + b \cdot \text{ID})$:}
\begin{align*}
n_k &\leftarrow n_k + 1 \\
\delta &= \varepsilon - \mu_k \\
\mu_k &\leftarrow \mu_k + \delta / n_k \\
M_k &\leftarrow M_k + \delta \cdot (\varepsilon - \mu_k)
\end{align*}

\textbf{Variance estimate:}
\[
\sigma_k^2 = \frac{M_k}{n_k - 1}
\]

\begin{itemize}
    \item Numerically stable --- avoids catastrophic cancellation in $\sum x^2 - (\sum x)^2 / n$.
    \item Streaming --- no need to store individual MT values.
    \item Per-pair --- captures that variance may differ across $(A, W)$ combinations.
\end{itemize}
\end{frame}

% ============================================================
\section{Tracker pipeline}

\begin{frame}{Tracker overview and motivation}
\textbf{Role of the tracker:} a Python process running on the Raspberry Pi that converts camera frames into 3D position estimates of the patient's device, transmitted to the Godot game via UDP at 30--60 Hz.

\textbf{Pipeline (every frame):}
\begin{enumerate}
    \item Capture grayscale frame from Pi camera.
    \item Undistort using fisheye intrinsics.
    \item Detect ArUco / AprilTag markers on the device.
    \item Refine corner positions to sub-pixel accuracy.
    \item Solve PnP per marker $\Rightarrow$ rotation $\mathbf{R}$ and translation $\mathbf{t}$.
    \item Apply marker-to-grip offsets, average across visible markers, transform into the world frame fixed at session start.
    \item Smooth temporally and transmit over UDP.
\end{enumerate}

\textbf{Why this matters for the study:} systematic tracker error propagates directly into game cursor position. Patient success rate becomes a function of \emph{tracker accuracy} as well as patient ability --- confounding the variable being studied.
\end{frame}

% ============================================================
\begin{frame}{Calibration mismatch on the Pi}
\textbf{Problem identified on inspection of inherited code:}

The original `main.py` loaded a hardcoded calibration file `calib_mono_faith.toml`. This file contained the intrinsics of a \emph{laptop webcam}, not the Pi camera.

\textbf{Consequence:} every \texttt{solvePnP} call on the Pi used the wrong intrinsic matrix $\mathbf{K}$ and the wrong distortion coefficients $\mathbf{D}$. Pose estimates were systematically biased --- though the system appeared to ``work'' because Godot received plausible-looking numbers.

\textbf{Why this is silent:} \texttt{solvePnP} does not validate that the supplied intrinsics match the camera that produced the image. Wrong intrinsics yield wrong but smooth pose estimates; the error is invisible without a ground-truth reference.

\textbf{Required fix:} produce a true calibration for the Pi camera and load it in place of the laptop's.
\end{frame}

% ============================================================
\begin{frame}{Camera calibration --- concepts}
\textbf{What calibration produces:}

A camera is described by two sets of parameters that map 3D world points to 2D image pixels.

\textbf{Intrinsic matrix:}
\[
\mathbf{K} = \begin{bmatrix} f_x & 0 & c_x \\ 0 & f_y & c_y \\ 0 & 0 & 1 \end{bmatrix}
\]
\begin{itemize}
    \item $f_x, f_y$ --- focal length in pixels (sensor-dependent).
    \item $c_x, c_y$ --- principal point (image centre, not always at the geometric centre of the sensor).
\end{itemize}

\textbf{Distortion coefficients:}
\begin{itemize}
    \item Pinhole model: $\mathbf{D} = [k_1, k_2, p_1, p_2, k_3]$ --- radial + tangential.
    \item Fisheye (Kannala-Brandt) model: $\mathbf{D} = [k_1, k_2, k_3, k_4]$ --- polynomial in incidence angle, accurate to large field-of-view.
\end{itemize}

The OV9281 has a $160^\circ$ fisheye lens. The pinhole model cannot fit this --- the fisheye model is required.
\end{frame}

% ============================================================
\begin{frame}{Fisheye calibration script}
\textbf{Implementation:} \texttt{pyscripts/calibrate\_camera.py}.

\textbf{Procedure:}
\begin{enumerate}
    \item Display a $9 \times 6$ chessboard pattern (square side $= 24.35$ mm, measured) at varying poses in front of the camera.
    \item For each frame, detect chessboard corners via \texttt{cv2.findChessboardCorners}.
    \item Auto-capture a frame when the corner positions have been stable for several frames (rejects motion blur).
    \item Collect $N \ge 20$ valid captures spanning different orientations and depths.
    \item Solve for $(\mathbf{K}, \mathbf{D})$ via \texttt{cv2.fisheye.calibrate} (Kannala-Brandt model).
    \item Save to \texttt{camera\_calib.toml}.
\end{enumerate}

\textbf{Multi-pose verification (post-calibration):} for $N$ held-still poses of the calibration board, report
\begin{itemize}
    \item \textbf{Accuracy} --- measured square edge length vs the true $24.35$ mm (tilt-invariant).
    \item \textbf{Precision} --- standard deviation of position across $N$ frames at the same pose.
    \item \textbf{Fit} --- mean reprojection error (px) of detected corners back to predicted positions.
\end{itemize}
\end{frame}

% ============================================================
\begin{frame}{Pipeline corrections in \texttt{main.py}}
After producing a valid Pi-camera calibration, three downstream bugs in the existing pipeline were corrected:

\begin{enumerate}
    \item \textbf{Calibration file selection moved to \texttt{settings.json}.} A single \texttt{calibration\_file} key replaces the hardcoded filename; the same code now loads the correct file on any machine.
    \item \textbf{Removed \texttt{cv2.flip(frame, 1)}.} A horizontal flip was being applied to every frame before detection. The flip mirrored the image across $x$, which made the detected marker centroids appear at $(W - u, v)$ instead of $(u, v)$. Since the calibration's principal point $c_x$ was fitted on un-flipped frames, the flip introduced a systematic $x$-axis offset in every pose estimate.
    \item \textbf{Removed \texttt{libcamera.Transform(vflip=1)}.} The picamera2 configuration was vertically flipping the sensor input. The OV9281 is mounted right-side up; the flip was a left-over from a different mounting orientation.
\end{enumerate}

\textbf{Resolution alignment:} \texttt{Config.FRAME\_SIZE} was \texttt{(1200, 800)}, but the calibration was performed at the OV9281 native \texttt{(1280, 800)}. Aligned to \texttt{(1280, 800)} so that downstream intrinsics match the captured frame dimensions.
\end{frame}

% ============================================================
\begin{frame}{Sensor-to-screen mapping --- 2D affine fit}
\textbf{Problem:} the tracker reports the device's 3D position in metres relative to a marker-anchored frame. The Godot game needs a 2D screen position in pixels. The mapping must absorb the camera-to-screen tilt and any axis misalignment between the workspace plane and the screen plane.

\textbf{Model (6-parameter affine transform):}
\begin{align*}
s_x &= a_{11} \, r_x + a_{12} \, r_z + a_{13} \\
s_y &= a_{21} \, r_x + a_{22} \, r_z + a_{23}
\end{align*}
where $(r_x, r_z)$ are the tracker's planar coordinates (metres) and $(s_x, s_y)$ are the screen position (pixels). The six coefficients capture scale, rotation, shear and translation.

\textbf{Calibration:} the patient touches four screen corners (top-left, top-right, bottom-left, bottom-right) with the device. Recorded as four $(r_x, r_z) \to (s_x, s_y)$ pairs. This yields $8$ equations for $6$ unknowns --- an over-determined system.

\textbf{Solved via ordinary least squares} (normal equations on a $3 \times 3$ system):
\[
\mathbf{N}^\top \mathbf{N} \, \boldsymbol{\theta} = \mathbf{N}^\top \mathbf{b}
\]
solved separately for each axis. The least-squares fit averages noise across the four corner samples rather than committing to a unique three-point solution.
\end{frame}

% ============================================================
\begin{frame}{Marker offsets from the CAD model}
\textbf{Problem:} when several markers are visible simultaneously, each one independently estimates the grip-point position via its own pose plus an offset vector. If the offsets are wrong, the per-marker estimates disagree and the centroid jumps when markers go in or out of view.

\textbf{Offset vector definition:} for each marker, an offset $\mathbf{o}$ expressed in the marker's \emph{own coordinate frame} (with $+\mathbf{X}$ = printed-right, $+\mathbf{Y}$ = printed-up, $+\mathbf{Z}$ = outward normal of marker face). The grip position in the camera frame is then
\[
\mathbf{p}_{\text{grip}} = \mathbf{R}_{\text{marker}} \, \mathbf{o} + \mathbf{t}_{\text{marker}}
\]

\textbf{Derivation procedure (Fusion 360):}
\begin{enumerate}
    \item Place a construction point at the grip location on the handle (centre of the cylindrical grip section).
    \item For each marker face, place a construction point at the centre of the mounting surface (intersection of diagonal construction lines).
    \item Read world coordinates of each point from Fusion's \texttt{Inspect} tool.
    \item For each marker, identify the face's outward normal $\hat{\mathbf{n}}$ in world coordinates (e.g.\ $+\mathbf{X}$ for the front face).
    \item Establish the marker's local frame in world coordinates: $\hat{\mathbf{Y}}_m = \text{device-up}$, $\hat{\mathbf{Z}}_m = \hat{\mathbf{n}}$, $\hat{\mathbf{X}}_m = \hat{\mathbf{Y}}_m \times \hat{\mathbf{Z}}_m$.
    \item Compute the world-frame displacement $\mathbf{d} = \mathbf{p}_{\text{grip}}^{\text{world}} - \mathbf{p}_{\text{face}}^{\text{world}}$ and project onto the local marker axes via dot products to obtain $\mathbf{o}$.
\end{enumerate}
\end{frame}

% ============================================================
\begin{frame}{Marker offsets --- deployed values}
\textbf{Marker orientation convention} (verified per marker via the live debug overlay):

Each side marker is glued with its printed $+\mathbf{Y}$ axis aligned to the device's upward direction. The printed $+\mathbf{Z}$ axis points outward from the marker face.

\textbf{Computed offsets (metres, marker-local frame):}
\[
\begin{array}{lll}
\text{Marker 12 (front face)} & \mathbf{o}_{12} & = (\phantom{-}0.001,\ \phantom{-}0.046,\ -0.059) \\
\text{Marker 14 (left face)}  & \mathbf{o}_{14} & = (-0.125,\ \phantom{-}0.045,\ -0.054) \\
\text{Marker 20 (right face)} & \mathbf{o}_{20} & = (\phantom{-}0.125,\ \phantom{-}0.045,\ -0.054)
\end{array}
\]

\textbf{Sanity checks:}
\begin{itemize}
    \item All three offsets have the same $+\mathbf{Y}$ component ($\sim 0.045$ m), consistent with the grip lying $\sim 4.5$ cm above each side-face centre.
    \item Marker 14 and Marker 20 offsets are mirror images about $\mathbf{X} = 0$, consistent with the device's left-right symmetry.
    \item All $\mathbf{Z}$ components are negative ($\sim -0.05$ m), reflecting that the grip is behind each outward-facing marker.
\end{itemize}
\end{frame}

% ============================================================
\section{Smoothing filters}

\begin{frame}{Motivation --- filtering the 3D position stream}
\textbf{Source of noise:} \texttt{solvePnP} produces a slightly different pose every frame, even when the marker is physically still. Two contributions:
\begin{itemize}
    \item Sub-pixel corner refinement noise ($\sigma \approx 0.1$--$0.3$ px per corner).
    \item Numerical conditioning of the PnP optimiser.
\end{itemize}

Translated through the rigid-body transform and the marker offsets, this produces a centroid that wiggles by a few millimetres frame-to-frame. The Godot cursor inherits the wiggle.

\textbf{Filter design constraints:}
\begin{itemize}
    \item Heavy smoothing when the patient holds still (suppress jitter).
    \item Minimal lag when the patient reaches (no perceived input delay).
    \item Computable in real time on a Raspberry Pi 5.
\end{itemize}

Three filters were added as opt-in alternatives, selected via \texttt{settings.json[``filter\_type'']}.
\end{frame}

% ============================================================
\begin{frame}{Exponential moving average (EMA)}
The original baseline filter. One smoothing parameter $\alpha \in [0, 1]$.

\textbf{Update rule:}
\[
\mathbf{y}_n = \alpha \cdot \mathbf{x}_n + (1 - \alpha) \cdot \mathbf{y}_{n-1}
\]

\textbf{Behaviour:}
\begin{itemize}
    \item $\alpha = 1$: output equals input (no smoothing).
    \item $\alpha = 0$: output equals previous output (input ignored).
    \item $\alpha = 0.4$ (deployed value): moderate smoothing, finite memory of past samples that decays geometrically.
\end{itemize}

\textbf{Limitations:}
\begin{itemize}
    \item Single fixed parameter $\Rightarrow$ cannot smooth heavily during a hold while remaining responsive during a reach. Whichever value is chosen, one regime suffers.
    \item No physical model of motion: treats velocity and acceleration as zero.
\end{itemize}
\end{frame}

% ============================================================
\begin{frame}{Kalman filter --- constant-velocity 3D}
A 6-dimensional state Kalman filter modelling position and velocity in 3D.

\textbf{State and measurement:}
\[
\mathbf{x} = \begin{bmatrix} x \\ y \\ z \\ \dot{x} \\ \dot{y} \\ \dot{z} \end{bmatrix} \in \mathbb{R}^6,
\qquad
\mathbf{z} = \begin{bmatrix} x \\ y \\ z \end{bmatrix} \in \mathbb{R}^3
\]

\textbf{Model matrices:}
\[
\mathbf{F}_n = \begin{bmatrix} \mathbf{I}_3 & \Delta t_n \cdot \mathbf{I}_3 \\ \mathbf{0} & \mathbf{I}_3 \end{bmatrix},
\qquad
\mathbf{H} = \begin{bmatrix} \mathbf{I}_3 & \mathbf{0} \end{bmatrix}
\]
\[
\mathbf{Q} = q \cdot \mathbf{I}_6 \ (\text{process noise}),
\qquad
\mathbf{R} = r \cdot \mathbf{I}_3 \ (\text{measurement noise})
\]

$\Delta t_n$ is measured per update from \texttt{time.monotonic()} so the filter remains correct under frame-rate jitter. Defaults: $q = 0.01$, $r = 0.05$.
\end{frame}

% ============================================================
\begin{frame}{Kalman filter --- predict and update}
\textbf{Predict step (every frame, using $\Delta t_n$):}
\begin{align*}
\hat{\mathbf{x}}_{n|n-1} &= \mathbf{F}_n \, \hat{\mathbf{x}}_{n-1|n-1} \\
\mathbf{P}_{n|n-1} &= \mathbf{F}_n \, \mathbf{P}_{n-1|n-1} \, \mathbf{F}_n^\top + \mathbf{Q}
\end{align*}

\textbf{Update step (using new measurement $\mathbf{z}_n$):}
\begin{align*}
\mathbf{S}_n &= \mathbf{H} \, \mathbf{P}_{n|n-1} \, \mathbf{H}^\top + \mathbf{R} \\
\mathbf{K}_n &= \mathbf{P}_{n|n-1} \, \mathbf{H}^\top \, \mathbf{S}_n^{-1} \\
\hat{\mathbf{x}}_{n|n} &= \hat{\mathbf{x}}_{n|n-1} + \mathbf{K}_n \bigl(\mathbf{z}_n - \mathbf{H} \, \hat{\mathbf{x}}_{n|n-1}\bigr) \\
\mathbf{P}_{n|n} &= (\mathbf{I} - \mathbf{K}_n \mathbf{H}) \, \mathbf{P}_{n|n-1}
\end{align*}

\textbf{Output:} the position components of $\hat{\mathbf{x}}_{n|n}$.

\textbf{Trade-off:} the constant-velocity assumption is violated at reach endpoints (when the patient stops at a target). The filter briefly overshoots after each stop because it has predicted continuing motion.
\end{frame}

% ============================================================
\begin{frame}{One Euro filter --- speed-adaptive smoothing}
\textbf{Core idea:} use the EMA update rule, but make the smoothing factor $\alpha$ change with the speed of the signal.
\begin{itemize}
    \item When the patient is \emph{still} (low speed): use small $\alpha$ $\Rightarrow$ heavy smoothing $\Rightarrow$ kills jitter.
    \item When the patient is \emph{reaching} (high speed): use large $\alpha$ $\Rightarrow$ light smoothing $\Rightarrow$ no perceived lag.
\end{itemize}

\textbf{Step 1 --- estimate the signal's speed.}
Compute the per-frame change $\dot{\mathbf{x}}_n = (\mathbf{x}_n - \mathbf{x}_{n-1}) / \Delta t$, then smooth it with a small EMA to reject noise:
\[
\hat{\dot{\mathbf{x}}}_n = \alpha_d \cdot \dot{\mathbf{x}}_n + (1 - \alpha_d) \cdot \hat{\dot{\mathbf{x}}}_{n-1}
\]

\textbf{Step 2 --- pick $\alpha$ from the estimated speed.}
\[
\alpha_c \;\text{grows with}\; \|\hat{\dot{\mathbf{x}}}_n\|
\quad\text{(slow $\Rightarrow$ small $\alpha$;\; fast $\Rightarrow$ large $\alpha$)}
\]

\textbf{Step 3 --- apply EMA to the position with the speed-dependent $\alpha$.}
\[
\hat{\mathbf{x}}_n = \alpha_c \cdot \mathbf{x}_n + (1 - \alpha_c) \cdot \hat{\mathbf{x}}_{n-1}
\]

The only difference from EMA: $\alpha_c$ is no longer fixed --- it adapts every frame based on how fast the cursor is moving. Reference: Casiez, Roussel \& Vogel, CHI 2012.
\end{frame}

% ============================================================
\begin{frame}{Corner stability filter --- gating \texttt{solvePnP}}
\textbf{Observation:} sub-pixel corner refinement noise produces frame-to-frame corner shifts of $\sim 0.1$--$0.3$ px even on a physically static marker. \texttt{solvePnP} amplifies this into millimetre-scale pose jitter via non-linear optimisation.

\textbf{Strategy:} if no corner has moved more than a threshold $\tau$ pixels since the last frame, do not re-run \texttt{solvePnP} --- reuse the previously computed pose.

\textbf{Stability check (every frame):}
\begin{enumerate}
    \item Detect markers and collect corners $\{\mathbf{c}_i^{(t)}\}$.
    \item Verify the same marker IDs and corner counts as the previous frame.
    \item Compute the maximum corner displacement:
    \[
    \Delta_{\max} = \max_i \, \bigl\| \mathbf{c}_i^{(t)} - \mathbf{c}_i^{(t-1)} \bigr\|
    \]
    \item If $\Delta_{\max} \le \tau$, mark as stable and reuse cached $(\mathbf{R}, \mathbf{t})$.
    \item Otherwise, run \texttt{solvePnP} and update the cache.
\end{enumerate}

\textbf{Effect:} eliminates jitter caused by optimiser variance on near-identical inputs. Deployed $\tau = 2$ px.
\end{frame}

\end{document}
